% Drawing 33/01-1 -- folder 33, batch 01, archive page 2 (PDF page 3).
% Transcription: feuille de croquis au crayon. En haut à droite « McL » et « $\times M'U$ », une lettre « A ». À gauche, un morphisme $X \to k$ et $k \to Y$ en équerre, puis $X' \to X$ avec une flèche oblique $X' \to k$ marquée \uncertain{« $\phi$ »}. En bas, deux esquisses $Y \to X$ au-dessus de $k$ et $Y' \to X'$ au-dessus de $k'$, une flèche $k \to k'$ barrée d'un petit cercle hachuré ; à droite un carré $X \leftarrow X$, $K$, $k'$, $k \leftarrow K$, un $X$ entouré.
% Kind: sheet of pencil sketches of arrow diagrams. Method 2 (schematic redraw from the facsimile). Status: partial.
% Notes: Redrawn with the sheet turned 180 degrees, the way the transcription reads it: then every letter is upright (McL, x M'U, A, X, k, Y, K, phi). Upper left: X over a vertical stroke down to k, and a horizontal stroke from k to Y (no arrowheads on the page). Below: X' -> X, an oblique arrow from X' down to k marked phi (the phi at the far left edge), and a vertical stroke from X down to that k. Upper right: McL, x M'U, and A with a faint short arrow after it. Lower half: Y -> X with a diagonal and a vertical stroke down to k; Y' -> X' with a curve and a vertical stroke down to k'; the stroke k -> k' crossed by a small hatched circle inside a faint larger ring, and a vertical arrow coming down across it; two long faint horizontal strokes under them at the left. Lower right: the square k', X (ringed) <- X, K, the arrow from the ringed X down to k, k -- K (no arrowhead on the page), the heavy vertical from K, the stroke from k' downwards, a faint arc over the top and a faint loop. Left out on purpose: the square Y <-h- Y', f, f', X <-g- X' in the middle right with its ring, "V'", "g(U')" and "U' != emptyset", which the transcription typesets as text (tikzcd and note) after the drawing mark; a faint dot with a short stroke left of A; the page number. Still differs: the faint pencil loops at the lower right are approximate. Redrawn 2026-11-09 (issue 128): to check again. Question settled by a machine pass on the facsimile, 2026-11-12 (page): turned 180 degrees every letter stands upright (McL, A, X, Y, k, K, phi, and the square Y, Y', X, X' with h, f, f', g), so the sheet was written that way up and scanned upside down; the vertical arrow coming down across the stroke k -> k' beside the hatched circle, with its head at the bottom, and the faint short stroke after A are both on the page and stay drawn. Not checked by a reader.
% Coordinates: 1 unit = 1 mm = 1/100 of the crop width (bbox in page fractions: [0.08, 0.06, 0.94, 0.88]).
% Redrawn by Claude (claude-opus-5-5), issue #73, 2026-11-09; not checked by a human; not compiled (no TeX in the session).
% Rights: reproduces part of the work -- behind RENDER_TRANSCRIPTS, never published by itself.
\documentclass[tikz,border=3pt]{standalone}
\usepackage{amsmath,amssymb}
\usetikzlibrary{patterns,arrows.meta}
\tikzset{pencil/.style={black!45}, struck/.style={black!35}, mark/.style={circle,draw,inner sep=1.3pt}, pt/.style={circle,fill,inner sep=0.9pt}, opt/.style={circle,draw,fill=white,inner sep=1.1pt}}
\begin{document}
\begin{tikzpicture}[x=1mm,y=1mm,line cap=round,line join=round,font=\scriptsize]
\node[inner sep=1pt] at (11.9,-12.7) {\small $X$};
\draw (11.2,-18) -- (10.8,-28.7);
\node[inner sep=1pt] at (9.7,-34.3) {\small $k$};
\draw (12.7,-34.3) -- (30.3,-34.8);
\node[inner sep=1pt] at (37,-34.7) {\small $Y$};
\node[inner sep=1pt] at (4.7,-49.3) {\small $X'$};
\draw[-{Stealth[length=2mm]}] (10.3,-49.9) -- (18,-50);
\node[inner sep=1pt] at (21.2,-49.7) {\small $X$};
\draw[-{Stealth[length=2mm]}] (5.7,-52.7) -- (16.7,-68.3);
\node[inner sep=1pt] at (1.2,-57.3) {\small $\phi$};
\draw (21,-54.7) -- (21.2,-66.7);
\node[inner sep=1pt] at (21.2,-69.7) {\small $k$};
\node[inner sep=1pt] at (79.7,-7) {\small McL};
\node[inner sep=1pt] at (81.3,-12.7) {\small $\times M'U$};
\node[inner sep=1pt] at (65.9,-18.7) {\small $A$};
\draw[-{Stealth[length=2mm]},pencil] (72.3,-18.9) -- (76,-18.9);
\node[inner sep=1pt] at (13.7,-85) {\small $Y$};
\draw[-{Stealth[length=2mm]}] (16.3,-85.3) -- (23.5,-85.5);
\node[inner sep=1pt] at (26.7,-85) {\small $X$};
\draw (15.2,-91.3) -- (22.1,-103.3);
\draw (26.7,-90.7) -- (26.5,-102.7);
\node[inner sep=1pt] at (27,-107) {\small $k$};
\node[inner sep=1pt] at (41.7,-84) {\small $Y'$};
\draw[-{Stealth[length=2mm]}] (45,-85.2) -- (51.7,-85.3);
\node[inner sep=1pt] at (53.7,-84) {\small $X'$};
\draw plot[smooth, tension=0.55] coordinates {(41.3,-88.7) (44.3,-96.7) (47.7,-107)};
\draw (52.7,-88.3) -- (50.5,-105.3);
\node[inner sep=1pt] at (52.7,-109) {\small $k'$};
\draw (30.8,-108.1) -- (48,-109.2);
\fill[pattern=north east lines] (38.9,-108) -- (38.7,-108.8) -- (38.2,-109.6) -- (37.5,-110) -- (36.7,-110.2) -- (35.8,-110) -- (35.1,-109.6) -- (34.6,-108.8) -- (34.5,-108) -- (34.6,-107.2) -- (35.1,-106.4) -- (35.8,-106) -- (36.7,-105.8) -- (37.5,-106) -- (38.2,-106.4) -- (38.7,-107.2) -- cycle;
\draw (36.7,-108) circle[radius=2.2];
\draw[pencil] (37.3,-106.5) ellipse[x radius=4.3, y radius=4.6];
\draw[-{Stealth[length=2mm]}] (44,-106.5) -- (43,-120);
\draw[pencil] (0,-116.3) -- (37.3,-115);
\draw[pencil] (0.3,-120.7) -- (36.3,-120.3);
\node[inner sep=1pt] at (62.3,-94.7) {\small $k'$};
\node[inner sep=1pt] at (70.7,-97.3) {\small $X$};
\draw[pencil] (70.5,-97) ellipse[x radius=3.3, y radius=5.3];
\draw[-{Stealth[length=2mm]}] (82.7,-97.2) -- (74,-97.3);
\node[inner sep=1pt] at (84.7,-97.7) {\small $X$};
\node[inner sep=1pt] at (88.3,-100) {\small $K$};
\draw[-{Stealth[length=2mm]}] (69.9,-103) -- (69.5,-113.3);
\node[inner sep=1pt] at (70,-116.3) {\small $k$};
\draw (71.3,-116.8) -- (80.3,-116.8);
\node[inner sep=1pt] at (84.7,-116) {\small $K$};
\draw[very thick] (85.5,-103.7) -- (84.8,-113);
\draw (63.9,-98.3) -- (66,-114);
\draw[pencil] plot[smooth, tension=0.55] coordinates {(71,-90.7) (78,-89.7) (85.7,-91.3)};
\draw[pencil] plot[smooth, tension=0.55] coordinates {(77.3,-93) (80,-104.7) (82.7,-114.7) (84.3,-108.7) (86.3,-94.7)};
\end{tikzpicture}
\end{document}
